医学视觉语言模型（VLM）通常用三个数字来概括：对图像退化的鲁棒性、对图像的依赖程度，以及跨医院迁移能力。人们往往把这三个数字当作模型自身的属性。我们表明，它们实际上是\textbf{模型--测量系统}的属性——既取决于模型行为，也取决于使该行为可被观测的测量设计；而且四个未被言明的前提决定了这些数字是否有意义：模型必须在真实图像上产生有变化的输出（\emph{响应性}）；反事实对照必须在同一度量空间中与零模型比较（\emph{零模型}）；召回率必须与工作点一并报告（\emph{工作点}）；回答契约与输出模式必须明确声明（\emph{契约覆盖}）。

我们在三个数据层上测试了六个开放权重胸片 VLM（两个胸片/医学模型、四个通用模型，来自四个家族；4-bit 为事后 NF4 量化，bf16/fp16 为发布精度）：域内 NIH ChestX-ray14（1,506 例研究 $\times$ 19 格合成退化网格）、跨医院 PadChest（1,500 例研究，存储分辨率加空白对照）以及报告生成 Open-i（947 例研究）。六模型矩阵共 \textbf{270,324} 次模型调用；计入仅跑散文层的第七个模型后为 \textbf{275,618} 次。所有调用使用同一套冻结的解码与预处理协议，并逐模型声明处理器几何。违反任一前提，读者本应得到的结论就会反转。若不做响应性检查，六个模型中有五个在 PadChest 上看似``对图像视而不见''；实际上它们对 1,500 张图像只给出 \textbf{1--20 种不同回答}，因此图像依赖不可测量，而且这一坍缩在 265 例疾病阳性研究中依然存在。若没有零模型，Qwen3-VL 在 Open-i 上看似高度依赖图像（Jaccard 0.09）；实际上它的空白回答是单条拒答模板，该对照测量的是``检测到图像缺失''。若不报告工作点，对比度退化看似使召回率\emph{提高}（+4.69\,\pp{}）；实际上阳性预测率同时上升 +4.06\,\pp{}，精确率下降，假阳性从 1,646 增至 1,905。若不声明输出模式，4-bit 量化看似使 MedGemma-1.5 的图像依赖降低到 2.21 分之一；实际上模型切换了解码体制（真实图像上的推理通道触发率由 0.13\% 变为 99.07\%），两个数字不在同一尺度上。

在这些控制之后，有两项正面结果成立。第一，退化作用于\textbf{两个可分离轴}——模型侧监测标记是否触发，以及退化对证据的\emph{损伤}程度——噪声落在危险象限：其损伤与分辨率损失相当（四个可读模型合并召回变化 $-0.128$ 对 $-0.169$），而标记触发率低 156 倍（0.46\% 对 71.98\%；该通道仅矩阵中的一个模型可读）。第二，退化敏感性的标签级排序\textbf{在 NIH 原始秩上具有模型不变性，但部分受患病率混杂}（四个可读模型两两 Spearman 0.66--0.86；对患病率做秩残差化后，偏 Spearman 在不同模型对间为 $-0.69$ 至 $+0.85$），这是候选设计规则，而不是已被证明的病灶规律。最后，我们报告一项负结果：复合风险信号只使发现级错误排序发生边际变化（AUROC 相对单独置信度的提升不超过 +0.034），且在 5\% 风险预算下没有模型能覆盖超过 9.3\% 的发现。

四项确证性实验在同 1,506 例研究与 \textbf{31,626} 个追加格子上补齐了原始设计留下的三个单检查点缺口（部分预算重跑处调用数更高）；schema 结果现已在第二个家族（Qwen3-VL）上复现，其余确证主张仍为单家族。仅用无患者重叠的训练集训练的 DenseNet-121 监督基线在七个病理标签上达到宏 AUROC \textbf{0.733}，而六个 VLM 为 0.50--0.63，说明``VLM 判别力在随机水平''是模型属性，而非任务不可学；同一 CNN 在 19 格网格上以固定 0.5 规则运行，也在对比度下移动了工作点（召回 +2.1\,\pp{}），因此反转 3 不是视觉语言模型特有的。在第二个医学模型上重做五套 schema 扫描，复现了空白图像的 schema 依赖性（五套 schema 的病理召回全部为 0.000），同时表明\emph{哪一种}失效形态由模型与契约共同决定。改写对照在 MedGemma-4B（$\kappa$ 0.62--0.64）与 Qwen3-VL（$\kappa$ 0.70--0.72）上都复现了真实图像结果，但空白图像答案对措辞的响应是模型特定的：p1 翻转 MedGemma 的空白答案，却不改变 Qwen3 的答案。第二个推理模型（Qwen3-VL-8B-Thinking）复现了工作点机制——阳性率由 0.250 降至 0.125，而宏 AUROC 始终平稳在 0.53--0.55——并暴露了原设计看不到的失效模式：在 5,120 token 预算下仅 \textbf{89.0\%} 的格子能给出可用答案（经 8,192 token 抽样后为 89.2\%），覆盖率本身随退化下降（92.6\% $\rightarrow$ 80.9\%），未完成的回复会撑满我们尝试过的每一个预算（上限 2,560、3,584、5,120、8,192 时，失败回复的中位长度分别为 9.8k、13.5k、19.3k、28.1k 字符）。两轴结构在两个模型上都成立（$|\Delta|$ 定义下可察觉性--损伤的 Spearman 均为 0.49，$n=6$，$p>0.26$），但象限归属并不通用：噪声对 MedGemma-1.5 的标记通道是``静默''的，对``是否终止''这一通道却是``响亮''的。静默失效是监测器的属性，而不是退化的属性。

我们不提出新模型。我们主张，依赖、鲁棒性与迁移应当作为\emph{可识别}的量来报告——这里的``可识别''取第~\ref{sec:preconditions}~节的审计意义，而非充要条件——并给出实现这一点的四项检查、相应工件与失效分类。阈值、配对、标签映射、聚类自助法与多重比较的敏感性分析均未改变四次反转；只在散文层运行的第七个模型作为``实测排除''报告；两个边际坍缩单元如实标为边际，而不作为证据使用。
